\begin{abstract}
Providers choose which benchmark results a release document prints, while independent evaluators score the same models on many more. Dating the instruments seems to settle what could have been reported, since a benchmark published after a model shipped could not be. The comparison fails on its own terms: the date that decides availability also fixes where a model sits in each benchmark's comparison pool, so benchmarks that postdate a release already trail those that predate it before any label enters. The failure is visible on a leaderboard we did not build, where fixed models keep identical scores across fourteen HELM Lite releases while every published win rate moves, and a companion leaderboard with an absolute headline moves nothing. Hand-coding every locatable release document shows the consequence: the coded gap rejects zero at $+11.4$ percentile points, yet removing every label returns $+14.1$ on the same releases. A study built this way turns pool composition into apparent concealment. What the record does establish is composition, reported benchmarks standing 8.7 points above omitted eligible ones within release, while no possession-grounded test detects withholding and the change-based drop gap keeps its predicted sign inside the permutation null. Carrying a fixed suite across releases would make later silence informative; adding labels to a moving pool does not.
\end{abstract}
